\SourceLeaf{238}%
gekrümmter Flächen folgen. Es ist jedoch von großer Bedeutung, auch diese Formel herzuleiten. Wir behalten die Zeichen aus Artikel 4 bei und setzen außerdem
\[
\begin{array}{lll}
\dfrac{\mathrm{d}\mathrm{d}x}{\mathrm{d}p^2}=\alpha,&\dfrac{\mathrm{d}\mathrm{d}x}{\mathrm{d}p\mathord{.}\mathrm{d}q}=\alpha',&\dfrac{\mathrm{d}\mathrm{d}x}{\mathrm{d}q^2}=\alpha''\\[6pt]
\dfrac{\mathrm{d}\mathrm{d}y}{\mathrm{d}p^2}=\beta,&\dfrac{\mathrm{d}\mathrm{d}y}{\mathrm{d}p\mathord{.}\mathrm{d}q}=\beta',&\dfrac{\mathrm{d}\mathrm{d}y}{\mathrm{d}q^2}=\beta''\\[6pt]
\dfrac{\mathrm{d}\mathrm{d}z}{\mathrm{d}p^2}=\gamma,&\dfrac{\mathrm{d}\mathrm{d}z}{\mathrm{d}p\mathord{.}\mathrm{d}q}=\gamma',&\dfrac{\mathrm{d}\mathrm{d}z}{\mathrm{d}q^2}=\gamma''
\end{array}
\]
Der Kürze halber setzen wir ferner
\[
\begin{aligned}
bc'-cb'&=A\\
ca'-ac'&=B\\
ab'-ba'&=C
\end{aligned}
\]
Zunächst bemerken wir, dass \(A\mathrm{d}x+B\mathrm{d}y+C\mathrm{d}z=0\) gilt, oder auch \(\mathrm{d}z=-\frac{A}{C}\mathrm{d}x-\frac{B}{C}\mathrm{d}y\). Soweit also \(z\) als Funktion von \(x,y\) betrachtet wird, ergibt sich
\[
\begin{aligned}
\frac{\mathrm{d}z}{\mathrm{d}x}&=t=-\frac{A}{C}\\
\frac{\mathrm{d}z}{\mathrm{d}y}&=u=-\frac{B}{C}
\end{aligned}
\]
Ferner leiten wir aus \(\mathrm{d}x=a\mathrm{d}p+a'\mathrm{d}q,\ \mathrm{d}y=b\mathrm{d}p+b'\mathrm{d}q\) ab, dass
\[
\begin{aligned}
C\mathrm{d}p&=b'\mathrm{d}x-a'\mathrm{d}y\\
C\mathrm{d}q&=-b\mathrm{d}x+a\mathrm{d}y
\end{aligned}
\]
Daraus erhalten wir die totalen Differentiale von \(t,u\)
\[
\begin{aligned}
C^3\mathrm{d}t&=(A\tfrac{\mathrm{d}C}{\mathrm{d}p}-C\tfrac{\mathrm{d}A}{\mathrm{d}p})(b'\mathrm{d}x-a'\mathrm{d}y)\\
&\quad +(C\tfrac{\mathrm{d}A}{\mathrm{d}q}-A\tfrac{\mathrm{d}C}{\mathrm{d}q})(b\mathrm{d}x-a\mathrm{d}y)\\
C^3\mathrm{d}u&=(B\tfrac{\mathrm{d}C}{\mathrm{d}p}-C\tfrac{\mathrm{d}B}{\mathrm{d}p})(b'\mathrm{d}x-a'\mathrm{d}y)\\
&\quad +(C\tfrac{\mathrm{d}B}{\mathrm{d}q}-B\tfrac{\mathrm{d}C}{\mathrm{d}q})(b\mathrm{d}x-a\mathrm{d}y)
\end{aligned}
\]
Wenn wir nun in diese Formeln
\[
\begin{aligned}
\frac{\mathrm{d}A}{\mathrm{d}p}&=c'\beta+b\gamma'-c\beta'-b'\gamma\\
\frac{\mathrm{d}A}{\mathrm{d}q}&=c'\beta'+b\gamma''-c\beta''-b'\gamma'\\
\frac{\mathrm{d}B}{\mathrm{d}p}&=a'\gamma+c\alpha'-a\gamma'-c'\alpha\\
\frac{\mathrm{d}B}{\mathrm{d}q}&=a'\gamma'+c\alpha''-a\gamma''-c'\alpha'\\
\frac{\mathrm{d}C}{\mathrm{d}p}&=b'\alpha+a\beta'-b\alpha'-a'\beta\\
\frac{\mathrm{d}C}{\mathrm{d}q}&=b'\alpha'+a\beta''-b\alpha''-a'\beta'
\end{aligned}
\]
